% ============================================================
\chapter{Lyapunov Stability}
\label{ch:lyapunov}
% ============================================================

\section{Introduction}

In 1892, a young Russian mathematician named Aleksandr Mikhailovich Lyapunov defended his doctoral thesis, \emph{The General Problem of the Stability of Motion}, at the University of Kharkov. The question he posed was disarmingly simple: if a system sits at equilibrium and we nudge it slightly, does it return, drift away, or do something else entirely? The answer, Lyapunov showed, does not require solving the differential equations explicitly. Instead, one can construct an auxiliary function---an ``energy-like'' quantity---and study how it evolves along trajectories.

This idea, now called the \emph{direct method} or \emph{second method} of Lyapunov, was revolutionary. It bypassed the hopeless task of finding closed-form solutions and replaced it with a conceptual argument: if you can find a function that decreases along every trajectory, the system must be converging. The catch, of course, is finding such a function---an art as much as a science, and one that remains at the frontier of research to this day. Alongside the direct method, Lyapunov also developed the \emph{indirect method}, which exploits linearisation near the equilibrium. Together, these two approaches form the backbone of stability theory.

\section{Definitions of Stability}

Consider the autonomous system $\dot{x} = f(x)$ where $f : U \to \R^n$ is
of class $C^1$ and $x^*$ is an equilibrium point: $f(x^*) = 0$. Without loss
of generality, we assume $x^* = 0$ (by translation).

\begin{definition}[Lyapunov stability]
The equilibrium point $x^* = 0$ is \textbf{stable in the sense of Lyapunov}
if for every $\epsilon > 0$, there exists $\delta > 0$ such that
\[
    \norm{x_0} < \delta \implies \norm{\varphi(t, x_0)} < \epsilon
    \quad \forall t \geq 0.
\]
\end{definition}

\begin{definition}[Asymptotic stability]
The equilibrium $0$ is \textbf{asymptotically stable} if it is Lyapunov
stable and there exists $\delta_0 > 0$ such that
\[
    \norm{x_0} < \delta_0 \implies \lim_{t \to +\infty} \varphi(t, x_0) = 0.
\]
\end{definition}

\begin{definition}[Global asymptotic stability]
The point $0$ is \textbf{globally asymptotically stable} (GAS) if it is
stable and $\lim_{t \to +\infty} \varphi(t, x_0) = 0$ for all $x_0 \in U$.
\end{definition}

\begin{definition}[Instability]
The point $0$ is \textbf{unstable} if it is not Lyapunov stable.
\end{definition}

\begin{definition}[Exponential stability]
The point $0$ is \textbf{exponentially stable} if there exist constants
$M > 0$, $\alpha > 0$, and $\delta > 0$ such that
\[
    \norm{x_0} < \delta \implies \norm{\varphi(t, x_0)} \leq M\norm{x_0}
    e^{-\alpha t} \quad \forall t \geq 0.
\]
\end{definition}

\begin{remark}
Exponential stability implies asymptotic stability, which implies Lyapunov
stability. The converses are false in general.
\end{remark}

\begin{intuition}
\begin{itemize}
    \item \textbf{Lyapunov stability}: a ball at the bottom of a bowl stays
    near the bottom if slightly perturbed.
    \item \textbf{Asymptotic stability}: the ball returns to the bottom due
    to friction.
    \item \textbf{Instability}: a ball atop a hill --- any perturbation
    drives it away.
\end{itemize}
\end{intuition}

\section{Basin of Attraction}

\begin{definition}[Basin of attraction]
The \textbf{basin of attraction} of an asymptotically stable equilibrium
$x^*$ is the set
\[
    \mathcal{B}(x^*) = \{x_0 \in U : \lim_{t \to +\infty} \varphi(t, x_0) = x^*\}.
\]
\end{definition}

\begin{proposition}
The basin of attraction $\mathcal{B}(x^*)$ is an open set invariant under
the flow, and its boundary $\partial\mathcal{B}(x^*)$ is also invariant.
\end{proposition}

\section{Indirect Method: Linearization}

\begin{theorem}[Stability by linearization --- Lyapunov indirect]
\label{thm:lyap-indirect}
Let $A = Df(0)$ be the Jacobian matrix of $f$ at the equilibrium $x^* = 0$.
\begin{enumerate}
    \item If all eigenvalues of $A$ have strictly negative real parts
    ($\text{Re}(\lambda_i) < 0$ for all $i$), then $0$ is
    \textbf{asymptotically stable}.
    \item If at least one eigenvalue of $A$ has strictly positive real part
    ($\text{Re}(\lambda_j) > 0$ for some $j$), then $0$ is \textbf{unstable}.
    \item If all eigenvalues satisfy $\text{Re}(\lambda_i) \leq 0$ with at
    least one eigenvalue on the imaginary axis, linearization is
    \textbf{inconclusive}: the direct method must be used.
\end{enumerate}
\end{theorem}

\begin{proof}[Proof sketch]
Write $f(x) = Ax + g(x)$ with $g(x) = o(\norm{x})$ near $0$. If all
eigenvalues of $A$ have negative real parts, there exists a positive
definite matrix $P$ such that $A^\top P + PA = -Q$ with $Q$ positive
definite (matrix Lyapunov equation). The function $V(x) = x^\top P x$ is
then a Lyapunov function for the linearized system, and the term
$g(x) = o(\norm{x})$ does not alter the conclusion in a sufficiently small
neighborhood of the origin.
\end{proof}

\begin{example}[Damped pendulum]
The linearized damped pendulum near $\theta = 0$ gives
\[
    A = \begin{pmatrix} 0 & 1 \\ -g/l & -c \end{pmatrix},
\]
with $c > 0$ the damping coefficient. The characteristic polynomial is
$\lambda^2 + c\lambda + g/l = 0$ with discriminant $\Delta = c^2 - 4g/l$.
Both eigenvalues have negative real parts (since $c > 0$ and $g/l > 0$),
so the equilibrium is asymptotically stable.
\end{example}

\section{Lyapunov's Direct Method}

\begin{definition}[Lyapunov function]
Let $V : U \to \R$ be a $C^1$ function defined on a neighborhood of
$x^* = 0$ with $V(0) = 0$. The \textbf{orbital derivative} (or Lie
derivative) of $V$ along the field $f$ is
\[
    \dot{V}(x) = \nabla V(x) \cdot f(x) = \sum_{i=1}^n \frac{\partial V}{\partial x_i}(x)\, f_i(x).
\]
We say that $V$ is:
\begin{itemize}
    \item \textbf{positive definite} if $V(x) > 0$ for all $x \neq 0$ in a
    neighborhood of $0$;
    \item \textbf{positive semi-definite} if $V(x) \geq 0$;
    \item \textbf{negative definite} if $V(x) < 0$ for all $x \neq 0$.
\end{itemize}
\end{definition}

\begin{theorem}[Lyapunov stability theorem]
\label{thm:lyap-stable}
If there exists a positive definite function $V$ such that $\dot{V}(x) \leq 0$
in a neighborhood of $0$, then $0$ is \textbf{Lyapunov stable}.
\end{theorem}

\begin{proof}
Let $\epsilon > 0$ be small enough that $\overline{B}(0, \epsilon) \subseteq U$.
Set $c = \min_{\norm{x} = \epsilon} V(x) > 0$ (since $V$ is positive definite
and continuous on the compact set $\{\norm{x} = \epsilon\}$). By continuity
of $V$ and $V(0) = 0$, there exists $\delta > 0$ such that $V(x) < c$ for
$\norm{x} < \delta$. If $\norm{x_0} < \delta$, then
$V(\varphi(t, x_0)) \leq V(x_0) < c$ for all $t \geq 0$ (since
$\dot{V} \leq 0$). Therefore $\varphi(t, x_0)$ cannot reach the sphere
$\norm{x} = \epsilon$, proving $\norm{\varphi(t, x_0)} < \epsilon$.
\end{proof}

\begin{theorem}[Lyapunov asymptotic stability theorem]
\label{thm:lyap-asymptotic}
If there exists a positive definite function $V$ such that $\dot{V}(x) < 0$
for all $x \neq 0$ in a neighborhood of $0$ (i.e., $\dot{V}$ is negative
definite), then $0$ is \textbf{asymptotically stable}.
\end{theorem}

\begin{theorem}[Barbashin-Krasovskii theorem]
\label{thm:barbashin}
If $V : \R^n \to \R$ is positive definite, radially unbounded
($V(x) \to +\infty$ as $\norm{x} \to +\infty$), and $\dot{V}(x) < 0$ for
all $x \neq 0$, then $0$ is \textbf{globally asymptotically stable}.
\end{theorem}

\begin{warning}[title=\textbf{Radially unbounded Lyapunov functions}]
For \emph{global} asymptotic stability, it is essential that $V$ be radially
unbounded. Without this condition, the trajectory could escape to infinity
while $V$ remains decreasing.
\end{warning}

\section{LaSalle's Invariance Principle}

When $\dot{V}$ is only negative semi-definite, the asymptotic stability
theorem does not apply directly. LaSalle's invariance principle often allows
one to conclude nonetheless.

\begin{theorem}[LaSalle's invariance principle]
\label{thm:lasalle}
Let $\Omega \subseteq U$ be a compact positively invariant set. Let
$V : \Omega \to \R$ be of class $C^1$ with $\dot{V}(x) \leq 0$ on $\Omega$.
Set $S = \{x \in \Omega : \dot{V}(x) = 0\}$ and let $M$ be the largest
invariant subset contained in $S$. Then
\[
    \omega(x_0) \subseteq M \qquad \forall x_0 \in \Omega.
\]
In particular, every trajectory starting in $\Omega$ converges to $M$.
\end{theorem}

\begin{proof}[Proof sketch]
Since $V$ is decreasing and bounded below on the compact set $\Omega$,
$V(\varphi(t, x_0))$ converges to a limit $c$. The $\omega$-limit set
$\omega(x_0)$ is nonempty (since $\Omega$ is compact), invariant, and
$V \equiv c$ on $\omega(x_0)$. Therefore $\dot{V} = 0$ on $\omega(x_0)$,
giving $\omega(x_0) \subseteq S$. By invariance of $\omega(x_0)$, we
obtain $\omega(x_0) \subseteq M$.
\end{proof}

\begin{example}[Damped pendulum --- LaSalle]
Consider the damped pendulum:
\[
    \dot{\theta} = \omega, \qquad \dot{\omega} = -\frac{g}{l}\sin\theta - c\omega,
    \quad c > 0.
\]
The energy function
$V(\theta, \omega) = \frac{1}{2}\omega^2 + \frac{g}{l}(1 - \cos\theta)$
satisfies
\[
    \dot{V} = -c\omega^2 \leq 0.
\]
The set $S = \{\dot{V} = 0\} = \{\omega = 0\}$. On $S$, we have
$\dot{\omega} = -\frac{g}{l}\sin\theta$, which vanishes only if
$\theta = k\pi$. The largest invariant set in $S$ is
$M = \{(0, 0), (\pi, 0)\}$. By LaSalle, every bounded trajectory converges
to $M$. Near the equilibrium $(0,0)$, we obtain asymptotic stability.
\end{example}

\section{Instability Theorems}

\begin{theorem}[Lyapunov instability theorem]
\label{thm:instability-lyap}
If there exists a $C^1$ function $V$ with $V(0) = 0$ such that
$\dot{V}(x) > 0$ in a neighborhood of $0$ (except at $0$), and if in every
neighborhood of $0$ there exists a point $x$ with $V(x) > 0$, then $0$ is
\textbf{unstable}.
\end{theorem}

\begin{theorem}[Chetaev's theorem]
Let $V : U \to \R$ be of class $C^1$ with $V(0) = 0$. Suppose there exists
an open set $G$ with $0 \in \partial G$ and:
\begin{enumerate}
    \item $V(x) > 0$ for all $x \in G$;
    \item $V(x) = 0$ on $\partial G \cap U$;
    \item $\dot{V}(x) > 0$ for all $x \in G$.
\end{enumerate}
Then $0$ is \textbf{unstable}.
\end{theorem}

\section{Construction of Lyapunov Functions}

The main difficulty of the direct method is \emph{finding} an appropriate
Lyapunov function. Here are several strategies.

\subsection{Quadratic Functions}

For the system $\dot{x} = Ax + g(x)$ with $g(x) = o(\norm{x})$, we seek
$V(x) = x^\top P x$ with $P$ a symmetric positive definite solution of the
matrix Lyapunov equation:
\[
    A^\top P + PA = -Q,
\]
where $Q$ is symmetric positive definite. This equation has a unique positive
definite solution $P$ if and only if $A$ is a Hurwitz matrix (all eigenvalues
with negative real parts).

\begin{keyformulas}[title=\textbf{Matrix Lyapunov equation}]
\[
    A^\top P + PA = -Q \qquad \Longrightarrow \qquad
    P = \int_0^{+\infty} e^{A^\top t}\, Q\, e^{At}\, dt.
\]
This integral converges if and only if $A$ is a Hurwitz matrix.
\end{keyformulas}

\subsection{Energy Functions}

For mechanical systems, the total energy (kinetic + potential) is often a good
candidate. If the system is dissipative, $\dot{V}$ will be negative or
semi-negative.

\subsection{Combinations and Estimates}

For more complex systems, one can construct $V$ by linear combination of
simple terms, or by the backstepping method in control theory.

\section{Stability of Periodic Orbits}

\begin{definition}[Orbital stability]
A periodic orbit $\Gamma$ is \textbf{orbitally stable} if for every
$\epsilon > 0$, there exists $\delta > 0$ such that
\[
    \mathrm{dist}(x_0, \Gamma) < \delta \implies
    \mathrm{dist}(\varphi(t, x_0), \Gamma) < \epsilon \quad \forall t \geq 0.
\]
\end{definition}

\begin{definition}[Asymptotic orbital stability]
The orbit $\Gamma$ is \textbf{asymptotically orbitally stable} if it is
orbitally stable and
$\mathrm{dist}(\varphi(t, x_0), \Gamma) \to 0$ as $t \to +\infty$
for $x_0$ sufficiently close to $\Gamma$.
\end{definition}

\begin{remark}
Orbital stability is more natural than Lyapunov stability for periodic orbits,
since even in the stable case, two nearby points on the orbit separate
linearly in time along the orbit direction.
\end{remark}

\section{Lyapunov Functions and Control}

\begin{remark}
In control theory, Lyapunov functions play a central role. For a controlled
system $\dot{x} = f(x, u)$, one seeks a state feedback $u = k(x)$ such that
a Lyapunov function $V$ satisfies $\dot{V} < 0$ in closed loop. This is the
principle of \textbf{Control Lyapunov Functions} (CLF).
\end{remark}

\section{Exercises}

\begin{exercise}[Stability by linearization]
Study the stability of the equilibrium points of the system
\[
    \dot{x} = y + x(1 - x^2 - y^2), \qquad \dot{y} = -x + y(1 - x^2 - y^2).
\]
\end{exercise}

\begin{exercise}[Lyapunov function construction]
For the system $\dot{x} = -x + 2x^2y$, $\dot{y} = -y$, show that
$V(x,y) = x^2 + y^2$ is a Lyapunov function in a suitable neighborhood
of the origin and conclude regarding stability.
\end{exercise}

\begin{exercise}[LaSalle]
Apply LaSalle's invariance principle to the system
\[
    \dot{x} = y, \qquad \dot{y} = -x^3 - y^3,
\]
with the function $V(x,y) = \frac{x^4}{4} + \frac{y^2}{2}$.
\end{exercise}

\begin{exercise}[Instability via Chetaev]
Show that the origin is unstable for the system
\[
    \dot{x} = y + x^3, \qquad \dot{y} = x + y^3,
\]
using the function $V(x,y) = xy$.
\end{exercise}

\begin{exercise}[Matrix Lyapunov equation]
Solve the equation $A^\top P + PA = -I$ for
$A = \begin{pmatrix} -1 & 0 \\ 1 & -2 \end{pmatrix}$.
Verify that $P$ is positive definite.
\end{exercise}

\begin{exercise}[Global stability]
Show that the origin is globally asymptotically stable for
$\dot{x} = -x^3$, $\dot{y} = -y$ using
$V(x,y) = \frac{x^4}{4} + \frac{y^2}{2}$.
\end{exercise}
